\documentclass{article}
\usepackage{amsmath,amssymb,xurl}
\usepackage[hidelinks]{hyperref}
% Shared commands for notes in docs/paper-gaps/.

\DeclareUnicodeCharacter{2080}{\ifmmode{}_{0}\else\textsubscript{0}\fi}
\DeclareUnicodeCharacter{2081}{\ifmmode{}_{1}\else\textsubscript{1}\fi}
\DeclareUnicodeCharacter{2082}{\ifmmode{}_{2}\else\textsubscript{2}\fi}
\DeclareUnicodeCharacter{2099}{\ifmmode{}_{n}\else\textsubscript{n}\fi}

\newcommand{\Lean}{\textsc{Lean}}
\ExplSyntaxOn
\NewDocumentCommand{\leanid}{m}{
  \ifmmode
    \textsf{\detokenize{#1}}
  \else
    \group_begin:
    \urlstyle{sf}
    \tl_set:Nn \l_tmpa_tl {#1}
    \tl_replace_all:Nnn \l_tmpa_tl {\_} {_}
    \exp_args:NV \path \l_tmpa_tl
    \group_end:
  \fi
}
\ExplSyntaxOff
\providecommand{\tr}{\operatorname{tr}}
\newcommand{\tnleanIssueBase}{https://github.com/LionSR/TNLean/issues/}
\newcommand{\tnleanPrBase}{https://github.com/LionSR/TNLean/pull/}
\newcommand{\tnleanDocBase}{https://sirui-lu.com/TNLean/docs/}
\newcommand{\ghissue}[1]{\href{\tnleanIssueBase#1}{GitHub issue \##1}}
\newcommand{\ghpr}[1]{\href{\tnleanPrBase#1}{PR \##1}}
\newcommand{\doclink}[1]{\href{\tnleanDocBase#1.html}{\path{#1}}}

% Dirac notation and the identity operator, as in blueprint/src/macros/common.tex.
% \providecommand keeps note-local definitions authoritative.
\usepackage{dsfont}
\providecommand{\ket}[1]{|#1\rangle}
\providecommand{\bra}[1]{\langle #1|}
\providecommand{\braket}[2]{\langle #1 | #2 \rangle}
\providecommand{\ketbra}[2]{|#1\rangle\!\langle #2|}
\providecommand{\Id}{\mathds{1}}
\providecommand{\one}{\mathds{1}}
\providecommand{\C}{\mathbb{C}}
\providecommand{\MN}[1]{M_{#1}(\C)}
\providecommand{\MD}{M_D(\C)}

% Verdict marker: \gapnote{<kind>}{<status>}. Kind is the policy
% classification (clarification, local-correction, scope-restriction,
% unfaithful, false-source, open-gap); status is open, wip (elimination
% underway), resolved, or historical. Machine-read by the site index;
% prints a verdict line.
\newcommand{\gapnote}[2]{\par\noindent\textbf{Verdict:} #1 (#2).\par}

\title{Native Dual Detours and Physical Bond Avoidance}
\author{}
\date{}
\begin{document}
\maketitle
\gapnote{scope-restriction}{open}

Schuch, Cirac, and P\'erez-Garc\'ia,
\href{https://arxiv.org/abs/1001.3807v3}{arXiv:1001.3807v3},
Definition 6.13 and Lemma 6.14, source
\path{Papers/1001.3807/paper_v3.tex}, lines 2180--2214, specify a native
bond insertion by a dual string and deform it away from an observation.

For a physical vertex set $R$, compute its forbidden dual plaquettes as
\begin{align}
    B(R)=\{(x,y):\{(x,y),(x+1,y),(x,y+1),(x+1,y+1)\}\cap R\ne\varnothing\}.
    \label{eq:physical_avoidance_forbidden}
\end{align}
The definition is \leanid{TNLean.PEPS.torusIncidentDualPlaquettes}.
Assume both torus periods are at least three, the actual closed-cell
realization of $B(R)$ is simply connected, and the endpoints of a given
four-labelled dual path $p$ are outside $B(R)$. The theorem
\leanid{TNLean.PEPS.exists_torusDualPath_zero_crossings_on_incident}
constructs a path $q$ with the same endpoints and lifted displacement such
that its signed crossing count is zero on every horizontal or downward
vertical physical bond with at least one endpoint in $R$.

The detour is derived from the exterior-collar theorem for $B(R)$.
Its visited plaquettes have all four corners outside $R$. Each native dual
step crosses a bond between two corners of one of its endpoint plaquettes,
so it crosses no incident physical bond. Summing these step counts proves
the assertion, including repeated or negatively oriented crossings.
Consequently \leanid{TNLean.PEPS.exists_torusDualPath_trivial_on_incident}
constructs such a path whose actual group insertions are identity on every
incident physical bond, for any group and insertion element. The group
need not be abelian; all insertion powers have the same base.

The simple connectedness hypothesis is on the explicitly computed dual
corner collar $B(R)$. It is not inferred merely from endpoint exclusion
or from simple connectedness of $R$. In particular no annular observation
with an enclosed endpoint is covered. The annular Wilson-loop obstruction
in \path{docs/paper-gaps/scp10_dual_flux_string_deformation.tex} continues to
rule out unrestricted endpoint-free vacuum-density assertions.

This closes the physical-incidence step under the stated collar geometry.
It does not prove homotopy completeness with forbidden square moves or the
open-pair vacuum-density equality. The latter still requires a complementary
physical contraction calculation; identity insertions on the observed side
alone do not supply it. These obligations remain under \ghissue{8703}.
The full-lattice Hamiltonian scope of \ghissue{8947} is unchanged.
\end{document}
